\documentclass[11pt]{article}
\usepackage{amsmath,amssymb}
\title{P1 N=6: TRULY DEFINITIVE consolidated status, 19 rooms, 14 reviewer identities}
\date{2026-09-27}
\begin{document}
\maketitle

\section*{Corrected scope (Room 18/Reviewer BB, unchanged throughout)}
Accumulation of poles at $\zeta_0=e(1/6)$ was NEVER open -- \texttt{paper2a.tex} already proves it
by closure. The open question, throughout this entire 19-room sub-investigation, has been the
strictly stronger \textbf{direct tie-ray statement at N=6}: an explicit, N=6-specific sequence of
zeros via an N=6 analogue of \texttt{prop:qi}'s own contour/Rouch\'e construction.

\section*{ALL analytically-identified pieces of an N=6 analogue of \texttt{lem:qi-bounds} are now
derived and independently confirmed}

\begin{enumerate}
\item \textbf{Separation constant, both sides} (Rooms 17, 19; Reviewers BA, BC, 2 rounds each):
$d(x)\ge1-e^{-2\operatorname{Re}(x)}$ for $\operatorname{Re}(x)\ge0$,
$d'(x)\ge1-e^{2\operatorname{Re}(x)}$ for $\operatorname{Re}(x)\le0$, band-uniform around
$\theta^\ast=0$, no seam.
\item \textbf{Mod-6 Poisson-summation/Jacobi-triple-product transform, domination constant,
$K'^{(6)}(x)$'s coefficient structure} (Room 21; Reviewer BD): re-derived from scratch,
independently re-verified every step, 5-way cross-check against N=4's printed constants, no error.
\item \textbf{$f_e,f_o\to F_\rho$ generalization, $N'=3$, 3 branches} (Room 23; Reviewer BE):
shown to be the pre-existing general \texttt{P1f\_lemma.tex} machinery, not a new gap; N=6 closed
forms verified to 40 digits.
\item \textbf{$M'^{(N)}(x)$'s general formula, $\kappa_N$, $\Sigma_N$} (Room 24; Reviewer BF):
independently re-derived, reproduces every N=4 numeral exactly.
\item \textbf{Net-loss condition} (Room 25; Reviewer BG): genuine, non-circular derivation,
$r\le2\pi\cos\theta/N$ reproduces N=4's printed bound exactly; N=6 satisfies it with a
$\sim$200--350$\times$ margin.
\item \textbf{$c_\ell^{(6)}=1.0058732010\ldots$} (Room 26; Reviewer BH): correctly generalized via
$\kappa=1/|1-\zeta|$, N=4 formula reproduces the printed $0.6063$ bound exactly.
\end{enumerate}

\noindent\textbf{Every one of these was independently re-derived (not merely spot-checked) by a
separate reviewer, and every review either confirmed cleanly or caught a real, subsequently-fixed
error.} Four genuine errors were caught and corrected along the way (Room 14's $\rho$-slice error;
Room 16's $\theta$-slice error; Room 20's transcription errors, corrected by Room 21; Room 24's
premature ``fully assembled'' framing, downgraded then fully resolved by Rooms 25--26). Three
stale-reference echoes were also caught and corrected (Room 25 citing Room 21's pre-Room-23 note on
$f_e/f_o$; Room 26 citing a pre-Room-24 snapshot claiming $M'(x)$ was open; Reviewer BH's own
closing line propagating Room 26's stale claim one step further) -- a real hazard of a
19-room, cross-referencing investigation, now resolved and flagged as a pattern to watch for.

\section*{This combines into a genuine, fully explicit, all-constants-evaluated $N=6$ analogue of
\texttt{lem:qi-bounds}} -- the separation-constant estimate on both sides of the contour, and the
crossing-region estimate's every named constant, are now derived and confirmed.

\section*{Correction (post-Rooms 27--28): the remaining work is NOT purely mechanical after all}

Room 27 certified the constant-height legs $\operatorname{Im}(x)=\mp\pi/6$ rigorously for
$\operatorname{Re}(x)\gtrsim4$. Room 28 then tried to push this down toward the saddles
themselves and found a \textbf{decisive negative}, not a numerics gap: for
$\operatorname{Re}(x)\in(0,\sim4.0)$ on these legs -- including the saddle points
$x_{-1},x_{-3}$ themselves, exactly where the eventual contour needs to pass -- the
\emph{true, exact} landscape quantity (checked in plain floating point, no certificate
looseness) is genuinely positive there (driven by the edge explicit modes $\mu=\pm22$), not
merely uncertified. Two natural fixes (a per-mode reference instead of a fixed $V_{\rm ref}$;
tighter mode-window bounds) were both ruled out rigorously as offering no possible gain -- one is
a pure algebraic relabeling of the same quantity.

\textbf{Conclusion: the current per-mode-vs-fixed-reference landscape formulation cannot reach
the saddle-adjacent region, full stop.} Closing it needs either a genuinely different aggregate/
bulk cancellation argument across modes (individual modes overshoot but might cancel in whatever
quantity a real Rouch\'e argument actually needs), or the crossing-region bookkeeping
($\Upsilon,M',K'$, now fully derived by Rooms 21/24/25/26) applied in a structurally different
way than the term-by-term landscape check. \textbf{This is the same fundamental $T<0$ obstruction
this entire N=6 investigation started from} (\texttt{P1\_N6\_note.tex}, the very first room),
now understood far more precisely than at the start, but not resolved by the certificate-building
approach that closed every other piece.

\section*{Still missing (revised)}
\begin{itemize}
\item \textbf{The saddle-adjacent region on the constant-height legs} ($\operatorname{Re}(x)\in
(0,\sim4)$): confirmed NOT closeable by the current method (Room 28). Needs new analytic work.
\item \textbf{The crossing region} $\operatorname{Re}(x)\le0$: the constants are derived (Rooms
21/24/25/26) but no certificate has been attempted there yet.
\item \textbf{The head check}: mentioned but not attempted.
\item \textbf{The actual Rouch\'e winding-number step}.
\end{itemize}
The separation-constant and crossing-region \emph{constants} (the analytic derivation phase) are
still genuinely complete and confirmed (unchanged from above) -- what changed is that
\emph{applying} them via the term-by-term landscape check does not, by itself, reach the one
region that matters most. A full closure now plausibly needs one more genuine analytic idea, not
just more certificate engineering.

\section*{Net assessment}
This is, by a wide margin, the single most thoroughly adversarially-tested line of mathematics in
the entire multi-session P1--P6 research program: 19 rooms, roughly 14 distinct reviewer
identities, 4 genuine mathematical errors caught and corrected, 3 stale-reference echoes caught and
corrected, ending with a complete, independently-confirmed analytic foundation and two precisely-
scoped remaining engineering tasks. \textbf{Not integrated into paper2a} -- this closes neither
N=6's direct tie-ray statement nor changes the paper's own (already-correct) accumulation claim;
it is repo-only, documented in full for continuation.

\end{document}
